\begin{abstract}
% Checkpoint A: Attention Getter / Motivation
Modern Security Operations Centers (SOCs) are overwhelmed by tens of thousands of security alerts daily, of which only a small fraction represent authentic malicious intrusions. Global empirical telemetry highlights this acute operational bottleneck: the IBM Cost of a Data Breach Report 2026 reveals an all-time high average breach cost of \$4.99 million (+12\% YoY) and an average lifecycle of 247 days from breach to containment, while Mandiant M-Trends 2026 reports a global median attacker dwell time of 14 days (surging to 122 days in cyber espionage incidents). Concurrently, national telemetry disclosures, such as Taiwan National Security Bureau's 2026 report of 2.63 million daily cyber probes targeting critical infrastructure (+6\% YoY), underscore the relentless external threat landscape. In this high-stakes environment, rapid and reliable triage is the critical determinant of organizational survival.

% Checkpoint B: But / Problem / Research Gap
However, massive volumes of false positives (FP) and benign triggers (BP) induce pervasive alert fatigue, risking catastrophic oversight of genuine attacks. In real-world enterprise telemetry (e.g., Microsoft GUIDE), genuine True Positives (TP) constitute only 20.6\% of incidents, whereas Benign Positives (49.0\%) and False Positives (30.4\%) dominate. Existing automated triage solutions suffer from fundamental shortcomings: standalone Machine Learning (ML) classifiers scale well across enterprise logs but lack attack kill-chain contextual semantics; conversely, recent empirical evaluations reveal that five mainstream tool-using LLM agent strategies suffer from a False Negative Rate (FNR) of at least 40.4\% on attack-related alerts, primarily due to premature termination when queries yield empty results. On the ExCyTIn-Bench benchmark, state-of-the-art LLMs achieve an investigation reward of only 0.606. Furthermore, opaque automation induces automation complacency among analysts, undermining human oversight.

% Checkpoint C: Cure / Proposed Solution
To overcome these foundational limitations, we propose the \textbf{AI SOC Analyst Agent}---a domain-specific, spec-driven AI system that synergistically orchestrates calibrated Machine Learning classification, MITRE ATT\&CK contextual reasoning, and bounded LLM agentic planning. The system delivers evidence-grounded, explainable triage for security incidents across TP, BP, and FP categories, equipped with a mathematically grounded human-in-the-loop escalation policy that knows precisely when an incident must be handed over to human analysts.

% Checkpoint D: Development / System Architecture
The agent employs an eight-module functional architecture. A Gradient Boosted Decision Tree (GBDT) model, trained on analyst-labeled enterprise telemetry, computes calibrated incident probabilities (TP/BP/FP) accompanied by TreeSHAP local feature explanations. Simultaneously, a vector database retrieves formal MITRE ATT\&CK STIX 2.1 techniques and NIST SP 800-61 Rev.~3 incident handling playbooks. An LLM planner conducts multi-step evidence gathering via structured SQL log query tools. Finally, an evidence fusion engine reconciles ML predictions, tactical kill-chain rules, and LLM reasoning. Grounded in the ``Assume Breach'' and ``Asymmetric Error Cost'' doctrines, the system systematically escalates incidents to human analysts whenever prediction confidence falls below threshold $\tau_{\text{conf}}$, conclusions conflict, investigative evidence is insufficient, or late-stage tactics (e.g., Exfiltration) are detected.

% Checkpoint E: Experiments / Evaluation Plan
We evaluate the proposed system on two benchmark suites: Microsoft GUIDE (encompassing $\sim$1 million real-world enterprise incidents across 6,100 customer organizations) for triage accuracy, and ExCyTIn-Bench (57 log tables, 8 authentic attack chains) for agentic multi-hop log investigation. Baselines include majority-class heuristics, detector-level historical false positive rules, standalone GBDT, and standalone LLM without knowledge retrieval. Primary evaluation metrics comprise Macro-F1, TP False Negative Rate, Human Escalation Rate, ATT\&CK mapping accuracy, report qualitative scoring, and execution cost (token usage and latency). Ablation studies isolate the quantitative contributions of knowledge retrieval, ML probability priors, and escalation logic.

% Checkpoint F: Expected Outcomes / Findings
As this project is in its active validation phase, we formally hypothesize: (1) integrating domain knowledge constraints and calibrated escalation rules will significantly reduce TP False Negative Rates compared to standalone ML and LLM baselines; (2) the operational trade-off curve between Miss Rate and Human Escalation Rate can be quantitatively mapped to provide actionable tuning for tiered SOC teams; and (3) structured ATT\&CK retrieval will demonstrably improve tactic mapping precision and report verifiability for Tier-1 analysts.
\end{abstract}

\begin{IEEEkeywords}
Security Operations Center (SOC), Alert Triage, LLM Agent, MITRE ATT\&CK, Explainable AI (TreeSHAP), Human-in-the-Loop, Spec-Driven Development.
\end{IEEEkeywords}
